\subsection{MAGMA, transport traces, and controlled factors}\label{sec:relevance}
The focused comparison adds the FP32 variable-size batched GEMM dispatcher
from MAGMA revision \texttt{c3de65f855dd} [12]. One composition follows
Boukaram et al.: each batch contains at most one block per output row,
and successive slots accumulate sequentially. A second composition batches
all block products into temporary storage, then reduces them by output row.
The latter tests whether avoiding sequential slots offsets its workspace
and reduction costs. Both cache dimensions and pointers, borrow device
values, and use the upstream GEMM dispatcher without kernel changes.
Setup validates dimensions before calling the dispatcher used by MAGMA's
\texttt{max\_nocheck} interface. Its normal launch and attribute-setting
overhead remains inside timing.

The replacement comparison crosses all 128 core configurations with matrix
seeds 2, 3, and 5, three fresh processes per seed, and batches of one or
eight products. Each process records twenty timing samples per method,
for 2,304 core processes. Eight-product samples divide elapsed time by eight.
They reuse the same input and output buffers and synchronize once per batch.
Methods and jobs are randomized, with all GPU processes run serially.
The comparator set includes the explicit CSR algorithms, grouped cuBLAS,
BSR8, applicable BSR32, and both MAGMA compositions. This campaign is separate
from the original NVIDIA-only campaigns. As noted above, only the 64 boundary
of the related hierarchical study's leaf sizes falls inside the supported
dimensions, so this comparison measures the compatible block product rather
than that study's larger leaves or its construction algorithm.

\paragraph{Repeatability.}
The headline and width plot use single-product timings. A configuration ratio
averages paired log ratios equally over three seeds and three processes per
seed. Table~\ref{tab:magma-repeatability} retains both timing modes. Its wins
use the three-process geometric mean for each configuration and seed.
Seed ranges average over processes, and process-label ranges average over
seeds. These descriptive ranges are not confidence intervals. Repeating both
factors avoids using timing trials as independent matrix samples.

\begin{table}[!htbp]
\centering\small
\caption{Repeated MAGMA-inclusive comparisons. Each width and batch contains
32 configurations, three seeds, and three processes per seed. Ratios above
one favor direct. Wins count 96 configuration and seed combinations.}
\label{tab:magma-repeatability}
\begin{tabular}{rrrrrrrrr}
\toprule
RHS & Batch & Best library & Seed range & Process range & Wins & Wins \(>5\%\) & Unanimous & Split\\
\midrule
8 & 1 & 1.060 & 1.057 to 1.063 & 1.057 to 1.064 & 81/96 & 23/32 & 24/32 & 7/32\\
16 & 1 & 1.045 & 1.043 to 1.049 & 1.043 to 1.047 & 72/96 & 22/32 & 24/32 & 0/32\\
32 & 1 & 1.088 & 1.084 to 1.090 & 1.087 to 1.088 & 68/96 & 20/32 & 20/32 & 4/32\\
64 & 1 & 1.398 & 1.396 to 1.401 & 1.397 to 1.399 & 96/96 & 32/32 & 32/32 & 0/32\\
8 & 8 & 1.052 & 1.048 to 1.060 & 1.048 to 1.060 & 71/96 & 21/32 & 21/32 & 5/32\\
16 & 8 & 1.032 & 1.031 to 1.035 & 1.031 to 1.033 & 72/96 & 21/32 & 24/32 & 0/32\\
32 & 8 & 1.079 & 1.076 to 1.082 & 1.078 to 1.079 & 63/96 & 20/32 & 20/32 & 3/32\\
64 & 8 & 1.395 & 1.391 to 1.397 & 1.394 to 1.397 & 96/96 & 32/32 & 32/32 & 0/32\\
\bottomrule
\end{tabular}

\end{table}

The overall ratios are \RelevanceCoreRatio{} for individual products and
\RelevanceQueuedRatio{} for batches of eight. The corresponding configuration
win counts, averaged across both seeds and processes, are
\RelevanceCoreWins{} and \RelevanceQueuedWins{} out of 128.

\paragraph{An application trace.}
A standalone periodic transport example solves \(u_t+u_x=0\) using
semi-Lagrangian discontinuous Galerkin projection in a Legendre modal basis [14].
Prescribed element orders repeat 7, 15, 31, and 63, giving block sizes
8, 16, 32, and 64. Each step translates the field by one quarter cell along
exact characteristics, then projects it onto the receiving element basis.
The self and upwind-neighbor overlap integrals define the two dense blocks
in each row. Their structure comes directly from the discretization.

All six cases reach the same physical time \(T=1/128\), with
\(\Delta t=h/4\). Element counts 256, 1,024, and 4,096 therefore take
8, 32, and 128 dependent products, respectively. Ensemble widths are 8 and 64.
Each case has three fresh processes and six trace samples per method.
A sample includes the entire alternating-buffer trace, an initial-state
reset, and output poisoning. All methods use the same assembled FP32 matrix.
Every final coefficient is compared with repeated double-precision sparse
products. The actual GPU output also passes a physical \(L^2\) comparison
with the analytically translated waves, with relative tolerance \(2\times10^{-3}\).
The executable must reject unchanged input before measuring any method.
Preparation requires more than a tenfold rejection margin on both coefficient
error criteria. A separate FP64 spatial-refinement study at fixed physical
time checks the discretization without the FP32 arithmetic floor.

This replaces the original near-identity forward-Euler trace. The old trace
remains historical evidence and does not contribute to current results.
The full overlap equations, refinement observations, and startup definition
are in the supplement. This remains a standalone transport component with
prescribed orders. It does not compare optimized DG-specific implementations.
It is nonetheless the closest thing to an application integration in this
paper. The operator comes from the discretization rather than from an
imposed partition, the products are dependent rather than repeated in
isolation, and the trace is validated against the analytic solution.

\begin{table}[!htbp]
\centering\small
\caption{Complete transport traces to \(T=1/128\). Ratios divide comparator time
by direct time and first average paired process ratios in log space.
The best-library column includes both MAGMA compositions and all tested
NVIDIA controls. Times are direct GPU milliseconds per trace.}
\label{tab:transport}
\begin{tabular}{rrrrrrr}
\toprule
Elements & RHS & Steps & Direct ms & MAGMA best & Best library & Wins\\
\midrule
256 & 8 & 8 & 0.108 & 2.290 & 0.840 & 1/3\\
256 & 64 & 8 & 0.223 & 2.069 & 1.065 & 3/3\\
1,024 & 8 & 32 & 0.681 & 3.934 & 1.186 & 3/3\\
1,024 & 64 & 32 & 2.392 & 2.085 & 1.244 & 3/3\\
4,096 & 8 & 128 & 12.697 & 2.971 & 1.220 & 3/3\\
4,096 & 64 & 128 & 50.824 & 1.833 & 1.439 & 3/3\\
\bottomrule
\end{tabular}

\end{table}

On the 4,096-element, RHS-64 transport case, the fastest-library-to-direct ratio is 1.438, with 3 of three direct wins. The ratio against the faster MAGMA composition is 1.823. The table retains smaller traces and narrow ensembles, including cases that favor another implementation. This is evidence for the stated dependent transport product, not a universal application speedup.


\paragraph{Separating workload factors.}
A crossed control varies block rows over \(128,512\), degree over
\(4,16,128\), block size over \(8,32\), and panel width over \(8,64\).
Each of the 24 configurations uses seed 11 and three fresh processes.
Degrees select nested prefixes of the same per-row column permutation,
and shape changes preserve that adjacency. The same library controls
are retained, with BSR32 added where compatible. Row-count changes also
change the working set, so they are not interpreted as isolated occupancy
measurements. The supplementary table gives every controlled combination.
These experiments distinguish the chosen factors without assigning a
single cause retrospectively to the covariance result.
Averaging equally over the other crossed factors, increasing degree from 4 to 128 changes the BSR8-to-direct ratio from 0.989 to 0.538 at RHS 8, and from 0.871 to 0.583 at RHS 64. Changing block size from 8 to 32 increases that ratio by factors of 1.708 and 3.206, respectively. The row-count contrast is much smaller at RHS 8 and unfavorable at RHS 64. Thus both degree and shape affect the performance boundary, and more block rows alone do not guarantee an advantage.

